\documentclass[a4paper,fontset=none]{ctexart}

\usepackage{anyfontsize}
\usepackage{amsmath,amssymb,mathrsfs,wasysym}
\usepackage{autobreak}
\allowdisplaybreaks
\usepackage[T1]{fontenc}
\usepackage{textcomp}
\usepackage{amsthm}
\usepackage[libertine,vvarbb]{newtxmath}
\usepackage[scr=rsfso,frak=euler,bb=ams]{mathalfa}
\usepackage{bm}
\usepackage{indentfirst}
\setlength{\parindent}{0em}
\usepackage{parskip}
\usepackage{geometry}
\geometry{top=3cm,bottom=3cm,left=2cm,right=2cm}
\usepackage{fontspec}
\setmainfont{Libertinus Serif}
\setsansfont{Libertinus Sans}
\setmonofont{JetBrains Mono}
\setCJKmainfont{SimSun}[BoldFont=Noto Sans CJK SC Medium]

\begin{document}

    \title{\textbf{电动力学作业}}
    \author{1900011604\ \ 张植竣\ \ 北京大学}
    \date{2021年4月5日}
    \maketitle

    \section*{第三章}

    \begin{itemize}

    \item[1.]
    由题意：
    \[
    \bm{B} = (0,0,B_z) = \nabla\times\bm{A} = (\partial_j A_k - \partial_k A_j, \partial_k A_i - \partial_i A_k, \partial_i A_j - \partial_j A_i)
    \]
    于是有：
    \[
    \partial_j A_k - \partial_k A_j = \partial_k A_i - \partial_i A_k = 0,\quad \partial_i A_j - \partial_j A_i = B_z
    \]
    若取：
    \[
    \partial_j A_k = \partial_k A_j = \partial_k A_i = \partial_i A_k = \partial_j A_i = 0,\quad \partial_i A_j = B_z
    \]
    又限定Coulomb规范$\nabla\cdot\bm{A}=\partial_i A_i+\partial_j A_j+\partial_k A_k=0$，取：
    \[
    \partial_i A_i = \partial_j A_j = \partial_k A_k = 0
    \]
    得磁矢势的一个取值：
    \[
    \bm{A}_1 = (0,B_z x,0)
    \]
    而若保留$\partial_j A_k=\partial_k A_j=\partial_k A_i=\partial_i A_k=\partial_j A_i=0$且$\partial_i A_j = B_z$，取规范：
    \[
    \partial_i A_i = \partial_j A_j = \partial_k A_k = 1
    \]
    得磁矢势的另一个取值：
    \[
    \bm{A}_2 = (x,B_z x + y,z)
    \]
    可见$\bm{A}_1$与$\bm{A}_2$相差一个矢量场$\bm{\varPsi}=(x,y,z)$。其旋度为：
    \[
    \nabla\times\bm{\varPsi} = (0,0,0) = \bm{0}
    \]
    即$\bm{A}_1$与$\bm{A}_2$之差是一个无旋矢量场。

    \item[3.(1)]
    由于$z$方向的平移对称性和$\theta$方向的旋转对称性，磁场应当只有$\theta$方向分量且大小只与$r$坐标有关。取Amp\'{e}re环路为圆环$r=R$，由Amp\'{e}re定律得：
    \[
    \left\{
        \begin{aligned}
            &2\pi RB_1=\mu I,\quad &z<0 \\
            &2\pi RB_2=\mu_0I,\quad &z>0
        \end{aligned}
    \right.
    \]
    边界条件为：
    \[
    \bm{n}\cdot(\bm{B}_2-\bm{B}_1) = B_{2z} - B_{1z} = 0
    \]
    解得：
    \[
    \bm{B} =
    \left\{
        \begin{aligned}
            &\frac{\mu I}{2\pi r}\hat{\theta},\quad &z<0 \\
            &\frac{\mu_0I}{2\pi r}\hat{\theta},\quad &z>0
        \end{aligned}
    \right.
    \]

    \item[3.(2)]
    介质表面没有自由电流，只有磁化电流：
    \[
    \bm{n}\times(\bm{M}_2-\bm{M}_1) = \bm{\alpha}_M
    \]
    而$\bm{H}=\dfrac{\bm{B}}{\mu_0}-\bm{M}$，得：
    \[
    \bm{M} = \frac{\bm{B}}{\mu_0}-\bm{H} =
    \left\{
        \begin{aligned}
            &\left(\frac{\mu}{\mu_0}-1\right)\frac{I}{2\pi r}\hat{\theta},\quad &z<0 \\
            &0,\quad &z>0
        \end{aligned}
    \right.
    \]
    故：
    \[
    \bm{\alpha}_M = -\hat{z}\times\left(\frac{\mu}{\mu_0}-1\right)\frac{I}{2\pi r}\hat{\theta} = \left(\frac{\mu}{\mu_0}-1\right)\frac{I}{2\pi r}\hat{r},\quad z=0
    \]
    在介质表面，由$\bm{M}$的Amp\'{e}re定律得磁化电流：
    \[
    I_M = \oint_L \bm{M}\cdot\mathrm{d}\bm{l} = \left(\frac{\mu}{\mu_0}-1\right)I,\quad z<0,r=0
    \]

    \item[4.(1)]
    求磁感应强度分布如下：

    根据对称性，$\bm{B}$的磁感应线应为$B$大小均匀的同心圆环，且$\bm{B}$在$x=0$的介质表面法向分量连续。
    
    又因为$\bm{B}=\mu\bm{H}$，$\displaystyle{\oint_L \bm{H}\cdot\mathrm{d}\bm{l}=I}$，说明在$x=0$的介质表面，$\bm{H}$将发生跃变。

    于是有：
    \[
    \oint_L \bm{H}\cdot\mathrm{d}\bm{l} = \frac{B}{\mu}\pi r + \frac{B}{\mu_0}\pi r = I,\quad B = \frac{\mu\mu_0 I}{(\mu+\mu_0)\pi r}
    \]
    即磁感应强度在真空和介质中均为：
    \[
    \bm{B} = \frac{\mu\mu_0 I}{(\mu+\mu_0)\pi r}\hat{\theta}
    \]

    \item[4.(2)]
    由$\bm{M}=\bm{B}\left(\dfrac{1}{\mu_0}-\dfrac{1}{\mu}\right)$得：
    \[
    \bm{M} =
    \left\{
        \begin{aligned}
            &\frac{(\mu-\mu_0)I}{(\mu+\mu_0)\pi r}\hat{\theta},\quad &x<0 \\
            &0,\quad &x>0
        \end{aligned}
    \right.
    \]
    则磁化电流$I_M$为：
    \[
    I_M = \oint_L \bm{M}\cdot\mathrm{d}\bm{l} = \pi r \cdot \frac{(\mu-\mu_0)I}{(\mu+\mu_0)\pi r} + \pi r \cdot 0 = \frac{(\mu-\mu_0)I}{(\mu+\mu_0)},\quad r=0
    \]

    \item[5.]
    由题意：磁场关于$z$轴中心对称，即$\dfrac{\partial B_\theta}{\partial\theta}=0$。

    而磁场是无散场：
    \[
    \nabla\cdot\bm{B} = \frac{1}{\rho}\cdot\frac{\partial}{\partial\rho}(\rho B_\rho) + \frac{1}{\rho}\cdot\frac{\partial B_\theta}{\partial\theta} + \frac{\partial B_z}{\partial z} = 0
    \]
    由于$B_z=B_0-C(z^2-\dfrac{1}{2}\rho^2)$，$\dfrac{\partial B_z}{\partial z}=-2Cz$。

    则$\dfrac{1}{\rho}\cdot\dfrac{\partial}{\partial\rho}(\rho B_\rho)=2Cz$，积分得：$B_\rho=Cz\rho+\dfrac{1}{\rho}C_1(z,\theta)$

    积分常数$C_1$与$\rho$无关，但可以与$z$和$\theta$有关。

    取$C_1(z,\theta)\equiv0$得：$B_\rho=Cz\rho$

    对于静磁场，应有$\nabla\times\bm{H}=0$，检验如下：

    由于$\bm{B}=\mu\bm{H}$，$\nabla\times\bm{H}=0$时，$\nabla\times\bm{B}=0$。

    由于$\dfrac{\partial B_\theta}{\partial\theta}=0$，可以取$B_\theta=0$，则$\bm{B}=Cz\rho\hat{\rho}+\left[B_0-C(z^2-\dfrac{1}{2}\rho^2)\right]\hat{z}$，其旋度为：
    \begin{align*}
    \nabla\times\bm{B}
    &= \frac{1}{\rho}(\partial_\theta B_z - \partial_z B_\theta)\hat{\rho} + (\partial_z B_\rho - \partial_\rho B_z)\hat{\theta} + \frac{1}{\rho}[\partial_\rho(\rho B_\theta) - \partial_\theta B_\rho]\hat{z} \\
    &= \frac{1}{\rho}(0 - 0)\hat{\rho} + (C\rho - C\rho)\hat{\theta} + \frac{1}{\rho}(0 - 0)\hat{z} \\
    &= 0
    \end{align*}
    综上，该处的$B_\rho=Cz\rho$。

    \item[7.]
    令$r=a$处的$\bm{A}=0$，在柱坐标系中，由于对称性，$\bm{A}$只有$z$方向分量且大小只与$r$有关，满足：
    \[
    \left\{
        \begin{aligned}
            &\nabla^2\bm{A}_1 = -\mu_0\bm{J},\quad &r<a \\
            &\nabla^2\bm{A}_2 = 0,\quad &r>a
        \end{aligned}
    \right.
    \]
    边界条件为：
    \[
    |\bm{A}_1|<\infty,\quad r=0
    \]
    \[
    \bm{A}=0,\quad \hat{r}\times\left(\frac{1}{\mu}\nabla\times\bm{A}_2 - \frac{1}{\mu_0}\nabla\times\bm{A}_1\right)=0,\quad r=a
    \]
    由于$A_r=A_\theta=0$，方程可以化简为：
    \[
    \left\{
        \begin{aligned}
            &\frac{1}{r}\cdot\frac{\mathrm{d}}{\mathrm{d}r}\left(r\frac{\mathrm{d}A_1}{\mathrm{d}r}\right) = -\mu_0 J,\quad &r<a \\
            &\frac{1}{r}\cdot\frac{\mathrm{d}}{\mathrm{d}r}\left(r\frac{\mathrm{d}A_2}{\mathrm{d}r}\right) = 0,\quad &r>a
        \end{aligned}
    \right.
    \]
    对方程直接积分：
    \[
    A_1 = -\frac{\mu_0}{4}Jr^2 + C_1\ln r + C_2
    \]
    \[
    A_2 = C_3\ln r + C_4
    \]
    由静磁场的唯一性定理，该特解是唯一解。

    因为$r=0$时$|\bm{A}_1|<\infty$，$C_1=0$。

    $r=a$处的切向边界条件可以化简为：
    \[
    \mu\frac{\mathrm{d}A_1}{\mathrm{d}r} = \mu_0\frac{\mathrm{d}A_2}{\mathrm{d}r}
    \]
    可得$C_3=-\dfrac{1}{2}\mu Ja^2$
    
    由$r=a$时$\bm{A}=0$得：$C_2=\dfrac{1}{4}\mu_0 Ja^2$，$C_4=-\dfrac{1}{2}\mu Ja^2\ln a$。

    故$\bm{A}$的分布为：
    \[
    \bm{A} =
    \left\{
        \begin{aligned}
            &\frac{1}{4}\mu_0(a^2-r^2),\quad &r<a \\
            &\frac{1}{2}\mu a^2 J\ln\frac{a}{r},\quad &r>a
        \end{aligned}
    \right.
    \]

    \item[8.]
    由$\bm{H}=\dfrac{q_m\bm{r}}{4\pi\mu_0 r^3}$得：$\bm{B}=\mu_0\bm{H}=\dfrac{q_m\bm{r}}{4\pi r^3}$。

    由 Stokes 公式，$\bm{B}=\nabla\times\bm{A}$的积分形式为：
    \[
    \Phi = \iint_\Sigma \bm{B}\cdot\mathrm{d}\bm{S} = \iint_\Sigma (\nabla\times\bm{A})\cdot\mathrm{d}\bm{S} = \oint_{\partial\Sigma} \bm{A}\cdot\mathrm{d}r
    \]
    考虑一个半径为$R$、极角为$\theta$的球冠：
    \[
    \Phi = \iint_\Sigma \bm{B}\cdot\mathrm{d}\bm{S} = \frac{q_m}{4\pi R^2}R^2 \int_0^\theta \sin\theta\,\mathrm{d}\theta \int_0^{2\pi}\mathrm{d}\phi = \frac{q_m}{2}(1-\cos\theta)
    \]
    球冠的底面圆周作为$\bm{A}$的积分路径：
    \[
    \oint_{\partial\Sigma} \bm{A}\cdot\mathrm{d}r = 2\pi A_\phi r\sin\theta
    \]
    则有：
    \[
    \bm{A} = A_\phi\hat{\phi} = \frac{q_m(1-\cos\theta)}{4\pi r\sin\theta} = \frac{q_m}{4\pi r}\tan\frac{\theta}{2}\hat{\phi}
    \]
    可见在$r=0$和$\theta=\pi$处，$\bm{A}$具有奇异性。

    \item[11.]
    设磁化球半径为$R$，球内外的磁标势分别为$\varphi_1$、$\varphi_2$。

    对于铁磁体球，联立$\bm{B}_1=\mu\bm{H}_1+\mu_0\bm{M}_0$和$\bm{B}_1=\mu_0(\bm{H}_1+\bm{M}_1)$得：
    \[
        \bm{M}_1 = \frac{\mu-\mu_0}{\mu\mu_0}\bm{B}_1 + \frac{\mu_0}{\mu}\bm{M}_0,\quad r<R
    \]
    对于外部的线性介质，联立$\bm{B}_2=\mu\bm{H}_2$和$\bm{B}_2=\mu_0(\bm{H}_2+\bm{M}_2)$得：
    \[
        \bm{M}_2 = \frac{\mu'-\mu_0}{\mu'\mu_0}\bm{B}_2,\quad r>R
    \]
    由Maxwell方程$\nabla\cdot\bm{B}=0$及$\bm{M}_0$为常矢量可得球内外磁荷：
    \[
        \rho_1=\nabla\cdot\bm{M}_1=0,\quad
        \rho_2=\nabla\cdot\bm{M}_2=0
    \]
    则$\varphi_1$、$\varphi_2$满足Laplace方程$\nabla^2\varphi_1=\nabla^2\varphi_2=0$，球坐标下可以用Legendre多项式展开：
    \[
    \varphi = \sum_{n=0}^\infty (a_n r^n + b_n r^{-n-1})\mathrm{P}_n(\cos\theta)
    \]
    由于$r=0$时$\varphi_1$有限，$r\to\infty$时$\varphi_2$有限，可得：
    \[
    \left\{
        \begin{aligned}
        &\varphi_1 = \sum_{n=0}^\infty a_n r^n\mathrm{P}_n(\cos\theta),\quad &r<R \\
        &\varphi_2 = \sum_{n=0}^\infty b_n r^{-n-1}\mathrm{P}_n(\cos\theta),\quad &r>R
        \end{aligned}
    \right.
    \]
    当$r=R$时，有边界条件：（法向量$\bm{n}=\hat{r}$）
    \[
    \varphi_1 = \varphi_2 = 0,\quad \bm{n}\cdot(\bm{B}_2-\bm{B}_1) = 0
    \]
    $r<R$时，$B_1=\mu H+\mu_0M_0\cos\theta=\mu\dfrac{\partial\varphi_1}{\partial r}+\mu_0M_0\cos\theta=-\mu\sum\limits_{n=0}^\infty n a_n r^{n-1}\mathrm{P}_n(\cos\theta)+\mu_0M_0\cos\theta$

    $r>R$时，$B_2=\mu'H=\mu'\dfrac{\partial\varphi_2}{\partial r}=\mu'\sum\limits_{n=0}^\infty (n+1)b_n r^{-n-2}\mathrm{P}_n(\cos\theta)$

    则有$r=R$处的边界条件：
    \[
    \left\{
        \begin{aligned}
        &-\mu\sum\limits_{n=0}^\infty n a_n R^{n-1}\mathrm{P}_n(\cos\theta)+\mu_0M_0\cos\theta = \mu'\sum\limits_{n=0}^\infty (n+1)b_n R^{-n-2}\mathrm{P}_n(\cos\theta) \\
        &\sum_{n=0}^\infty a_n R^n\mathrm{P}_n(\cos\theta) = \sum_{n=0}^\infty b_n R^{-n-1}\mathrm{P}_n(\cos\theta)
        \end{aligned}
    \right.
    \]
    比较$\mathrm{P}_n(\cos\theta)$的各阶系数得：
    \[
    \left\{
        \begin{aligned}
        &a_n = b_n = 0,\quad &n\neq1 \\
        &a_n = \frac{\mu_0 M_0}{2\mu'+\mu},\quad b_n = \frac{\mu_0 M_0 R^3}{2\mu'+\mu},\quad &n=1
        \end{aligned}
    \right.
    \]
    故有：
    \[
    \varphi_1 = \frac{\mu_0 M_0 r}{2\mu'+\mu}\cos\theta,\quad
    \bm{H}_1 = -\nabla\varphi_1 = -\frac{\mu_0\bm{M}_0}{2\mu'+\mu}
    \]
    \[
    \varphi_2 = \frac{\mu_0 M_0 R^3}{(2\mu'+\mu)r^2}\cos\theta,\quad
    \bm{H}_2 = -\nabla\varphi_2 = \frac{\mu_0 R^3}{2\mu'+\mu}\left[\frac{3(\bm{M}_0\cdot\bm{r})\bm{r}}{r^5}-\frac{\bm{M}_0}{r^3}\right]
    \]
    而磁感应强度：
    \[
        \bm{B}_1 = \mu\bm{H}_1 + \mu_0\bm{M}_0,\quad
        \bm{B}_2 = \mu'\bm{H}_2
    \]
    代入得：
    \[
    \bm{B} =
    \left\{
        \begin{aligned}
            &\frac{2\mu'\mu_0 \bm{M}_0}{2\mu'+\mu},\quad &r<R \\
            &\frac{\mu'\mu_0 R^3}{2\mu'+\mu}\left[\frac{3(\bm{M}_0\cdot\bm{r})\bm{r}}{r^5}-\frac{\bm{M}_0}{r^3}\right],\quad &r>R
        \end{aligned}
    \right.
    \]
    磁化电流$\bm{\alpha}_M$只分布在球面$r=R$处，与$\bm{M}$的切向分量相关：$\bm{\alpha}_M=\bm{n}\times(\bm{M}_2-\bm{M}_1)$。
    
    由于边界条件$\varphi_1=\varphi_2$蕴含了$r=R$处$\bm{n}\times(\bm{H}_2-\bm{H}_1)=0$，磁化电流为：
    \[
    \bm{\alpha}_M = \bm{n}\times(\bm{M}_2-\bm{M}_1) = \frac{1}{\mu_0}\bm{n}\times(\bm{B}_2-\bm{B}_1) = \frac{3\mu'\mu_0}{2\mu'+\mu}\bm{M}_0\sin\theta\hat{\phi},\quad r=R
    \]
    注意此处$\bm{B}_2$的第一项$\dfrac{\mu'\mu_0 R^3}{2\mu'+\mu}\left[\dfrac{3(\bm{M}_0\cdot\bm{r})\bm{r}}{r^5}\right]$是径向分量，对切向分量没有贡献。

    \item[14.(1)]
    设球内的电荷密度为$\rho$，则$\rho=\dfrac{3Q}{4\pi R_0^3}$。
    
    考虑一个半径为$r$，高度为$r\mathrm{d}\theta$，沿球半径方向的厚度为$\mathrm{d}r$的环，其体积为$\mathrm{d}V=2\pi r^2\sin\theta\mathrm{d}r\mathrm{d}\theta$。
    
    环内的电流为：
    \[
    I = \frac{\mathrm{d}q}{\mathrm{d}t} = \frac{\Delta q}{T} = \frac{\rho \cdot 2\pi r^2\sin\theta\mathrm{d}r\mathrm{d}\theta}{2\pi\omega^{-1}} = \rho\omega r^2\sin\theta\mathrm{d}r\mathrm{d}\theta
    \]
    环的磁矩为：
    \[
    \mathrm{d}\bm{m} = I\bm{S} = \rho\omega r^2\sin\theta\mathrm{d}r\mathrm{d}\theta \cdot \pi(r\sin\theta)^2\hat{z} = \pi\rho\bm{\omega}r^4\sin^2\theta\mathrm{d}r\mathrm{d}\theta
    \]
    其中$\bm{S}=\pi(r\sin\theta)^2\hat{z}$为电流环围成的圆面的有向面积。
    
    则球的总磁矩为：
    \[
    \bm{m} = \iiint\limits_V \mathrm{d}\bm{m} = \pi\rho\bm{\omega}\int_0^{R_0} r^4\,\mathrm{d}r \int_0^\pi \sin^3\theta\,\mathrm{d}\theta = \frac{4}{15}\pi\rho\bm{\omega}R_0^5
    \]
    代入$\rho=\dfrac{3Q}{4\pi R_0^3}$得：$\bm{m}=\dfrac{1}{5}Q\bm{\omega}R_0^2$

    \item[14.(2)]
    设球的质量密度为$\rho_m$，则$\rho_m=\dfrac{3m_0}{4\pi R_0^3}$。
    
    考虑与(1)中相同的一个圆环，其角动量为：
    \[
    \mathrm{d}\bm{L} = \rho_m\mathrm{d}V \cdot \bm{\omega}(r\sin\theta)^2 = 2\pi\rho_m\bm{\omega}r^4\sin^3\theta\mathrm{d}r\mathrm{d}\theta
    \]
    则球的总角动量为：
    \[
    \bm{L} = \iiint\limits_V \mathrm{d}\bm{L} = 2\pi\rho_m\bm{\omega}\int_0^{R_0} r^4\,\mathrm{d}r \int_0^\pi \sin^3\theta\,\mathrm{d}\theta = \frac{8}{15}\pi\rho_m\bm{\omega}R_0^5
    \]
    代入$\rho_m=\dfrac{3m_0}{4\pi R_0^3}$得：$\bm{L}=\dfrac{2}{5}m_0\bm{\omega}R_0^2$，两者的比值为$\dfrac{\bm{m}}{\bm{L}}=\dfrac{Q}{2m_0}$。

    \item[15.]
    设介质表面为$z=0$平面，$\bm{m}$位于介质表面上方$z=a$处，与界面法向$\hat{z}$的夹角为$\theta$。

    由于高磁导率介质的表面为等磁标势面，考虑使用磁镜像法处理。
    
    小永磁体$\bm{m}$的镜像$\bm{m}'$位于介质表面下方$z=-a$处，与界面法向$\hat{z}$的夹角为$-\theta$，与$\bm{m}$的夹角为$2\theta$。

    $\bm{m}'$在介质外部任意一点的磁标势与磁感应强度为：
    \[
        \varphi' = \frac{\bm{m}'\cdot\bm{r}}{4\pi r^3},\quad \bm{B}' = -\mu_0\nabla\varphi' = \frac{\mu_0}{4\pi}\left[\frac{3(\bm{m}'\cdot\bm{r})\bm{r}}{r^5}-\frac{\bm{m}'}{r^3}\right]
    \]
    则$\bm{m}'$在$\bm{m}$处（距离为$2a$，$\bm{m}'$与$\bm{r}$夹角为$-\theta$）产生的磁感应强度为：
    \[
        \bm{B}' = \frac{\mu_0}{32\pi a^3}(3m'\cos\theta\hat{z}-\bm{m}')
    \]
    注意到$m'=m$，得$\bm{m}'$的磁场对$\bm{m}$的作用能为：
    \[
        W = -\bm{m} \cdot \bm{B}' = -\frac{\mu_0 m^2}{32\pi a^3}\left(1+\cos^2\theta\right)
    \]
    作用在小永磁体上的磁场力$\bm{F}$为：
    \[
        \bm{F} = -\nabla W = -\frac{\partial W}{\partial z}\hat{z} = -\frac{3\mu_0 m^2}{64\pi a^4}\left(1+\cos^2\theta\right)\hat{z}
    \]

    \end{itemize}

    \subsection*{补充题}

    \begin{itemize}

    \item[2.]
    在柱坐标中处理，首先根据对称性：$\dfrac{\partial\bm{B}}{\partial z}=\dfrac{\partial\bm{B}}{\partial\theta}=0$。
    
    又因为电流沿$\hat{\theta}$方向，磁场方向只能沿$\hat{z}$和$\hat{\rho}$方向。

    而磁场是无散场：
    \[
    \nabla\cdot\bm{B} = \frac{1}{\rho}\cdot\frac{\partial}{\partial\rho}(\rho B_\rho) + \frac{1}{\rho}\cdot\frac{\partial B_\theta}{\partial\theta} + \frac{\partial B_z}{\partial z} = 0
    \]
    故$\dfrac{\partial}{\partial\rho}(\rho B_\rho)=0$，不妨假设$B_\rho=\dfrac{\partial B_\rho}{\partial\rho}(\rho B_\rho)=0$，内部磁场为沿$\hat{z}$方向的匀强磁场。

    无穷远处的边界条件为$\lim\limits_{\rho\to\infty}\bm{B}_2=\lim\limits_{\rho\to\infty}\bm{H}_2=0$，由于螺线管无穷长，外部磁场$\bm{B}_2=\mu_0\textbf{H}_2=0$。

    由Amp\'{e}re环路定理，在$\theta=0$的平面取一个长方形环路，其中两条边平行于$z$轴、长度为$l$，则有：
    \[
    \mu_0(nl)I = lB_1,\quad B_1 = \mu_0nI
    \]
    这符合$\nabla\cdot\bm{B}=0$，$\nabla\times\bm{H}=0$和全部边界条件，由惟一性定理得：
    \[
    \bm{B} =
    \left\{
        \begin{aligned}
            &\mu_0nI\hat{z} &r<R \\
            &0,\quad &r>R
        \end{aligned}
    \right.
    \]

    \item[5.]
    由于磁标势$\varphi$满足Laplace方程，可以给出$\varphi$的形式：
    \[
    \varphi = \sum_{n=0}^\infty (a_n r^n + b_n r^{-n-1})\mathrm{P}_n(\cos\theta)
    \]
    当$r=R$时，有边界条件：（此处$\bm{n}=\hat{r}$）
    \[
    \varphi_1 = \varphi_2 = 0,\quad \bm{n}\cdot(\bm{B}_2 - \bm{B}_1) = \mu_0\frac{\partial\varphi_2}{\partial r} - \mu\frac{\partial\varphi_1}{\partial r} = 0
    \]
    设外场$\bm{H}_0=H_0\hat{z}=H_0\hat{r}\cos\theta$，当$r\to\infty$时，有边界条件：（注意$\bm{H}=-\nabla\varphi$）
    \[
    \lim_{r\to\infty}\varphi_2 = -H_0 r\cos\theta
    \]
    于是有：
    \[
    \left\{
        \begin{aligned}
        &\varphi_1 = \sum_{n=0}^\infty a_n r^n\mathrm{P}_n(\cos\theta),\quad &r<R \\
        &\varphi_2 = -H_0 r\cos\theta + \sum_{n=0}^\infty b_n r^{-n-1}\mathrm{P}_n(\cos\theta),\quad &r>R
        \end{aligned}
    \right.
    \]
    则有$r=R$处的边界条件：
    \[
    \left\{
        \begin{aligned}
        &\mu\sum\limits_{n=0}^\infty n a_n R^{n-1}\mathrm{P}_n(\cos\theta) = -\mu_0H_0R\cos\theta + \mu_0\sum\limits_{n=0}^\infty (n+1)b_n R^{-n-2}\mathrm{P}_n(\cos\theta) \\
        &\sum_{n=0}^\infty a_n R^n\mathrm{P}_n(\cos\theta) = -H_0R\cos\theta + \sum_{n=0}^\infty b_n R^{-n-1}\mathrm{P}_n(\cos\theta)
        \end{aligned}
    \right.
    \]
    比较$\mathrm{P}_n(\cos\theta)$的各阶系数得：
    \[
    \left\{
        \begin{aligned}
        &a_n = b_n = 0,\quad &n\neq1 \\
        &a_n = -\frac{3\mu_0 H_0}{2\mu_0+\mu},\quad b_n = \frac{(\mu-\mu_0) H_0 R^3}{2\mu_0+\mu},\quad &n=1
        \end{aligned}
    \right.
    \]
    故有：
    \[
    \varphi_1 = -\frac{3\mu_0 H_0r}{2\mu_0+\mu}\cos\theta,\quad
    \bm{H}_1 = -\nabla\varphi_1 = \frac{3\mu_0\bm{H}_0}{2\mu_0+\mu}
    \]
    \[
    \varphi_2 = -H_0r\cos\theta + \frac{(\mu-\mu_0) H_0 R^3}{(2\mu_0+\mu)r^2}\cos\theta,\quad
    \bm{H}_2 = -\nabla\varphi_2 = \bm{H}_0 + \frac{(\mu-\mu_0)R^3}{2\mu_0+\mu}\left[\frac{3(\bm{H}_0\cdot\bm{r})\bm{r}}{r^5}-\frac{\bm{H}_0}{r^3}\right]
    \]
    而磁感应强度为：
    \[
    \bm{B}_1 = \mu\bm{H}_1 = \frac{3\mu\mu_0\bm{H}_0}{2\mu_0+\mu}
    \]
    \[
    \bm{B}_2 = \mu_0\bm{H}_2 = \mu_0\bm{H}_0 + \frac{\mu_0(\mu-\mu_0)R^3}{2\mu_0+\mu}\left[\frac{3(\bm{H}_0\cdot\bm{r})\bm{r}}{r^5}-\frac{\bm{H}_0}{r^3}\right]
    \]
    球体内$\bm{B}_1=\mu\bm{H}_1=\mu_0(\bm{H}_1+\bm{M})$，则磁化强度：
    \[
    \bm{M} = \frac{\bm{B}_1}{\mu_0} - \bm{H}_1 = \frac{3(\mu-\mu_0)}{2\mu_0+\mu}\bm{H}_0
    \]
    球体磁矩为：
    \[
    \bm{m} = \frac{4\pi R^3}{3}\bm{M} = \frac{4\pi R^3(\mu-\mu_0)}{2\mu_0+\mu}\bm{H_0}
    \]

    \end{itemize}

\end{document} 
